\section{Research Gap}
\label{sec:research_gap}

The reviewed literature establishes several important principles for exploration in reinforcement learning.
Novelty-based methods encourage state-space coverage and can guide agents toward sparse rewards, but they do not directly measure whether the agent is learning from the visited states.
Prediction-error curiosity provides a neural and representation-based form of surprise, but it can reward stochastic or unlearnable regions.
Impact-driven exploration encourages controllable change, but impact alone does not determine whether the resulting experience improves the agent's model or policy.
Learning-progress methods address the difference between surprise and improvement, but they require stable local estimates and an appropriate region structure in high-dimensional or continuous spaces.

A remaining gap is the need for an intrinsic-reward formulation that combines these complementary strengths while reducing their individual failure modes.
In particular, an exploration signal for deep reinforcement learning should satisfy four requirements.
First, it should operate in a learned representation so that it can be used with continuous observations and neural policies.
Second, it should value transitions that produce meaningful feature-space change, since exploration should involve behaviorally consequential action.
Third, it should estimate learning progress locally, because different parts of the environment may become predictable at different rates.
Fourth, it should suppress intrinsic reward in regions where prediction error remains high but learning progress remains low.

Existing methods address parts of this set of requirements but not the full combination.
ICM and RND provide scalable prediction-error curiosity, but their intrinsic rewards are primarily driven by error magnitude \cite{pathak-2017,burda-2018B}.
RIDE introduces a feature-space impact signal and count-based modulation, but it does not directly use local model-improvement estimates as the main criterion for productive exploration \cite{raileanu-2020}.
R-IAC uses region-local learning progress, but adapting this idea to deep reinforcement learning requires learned latent representations, stable scale handling, and integration with modern policy-gradient training \cite{baranes-2009}.
The noisy-TV literature further shows that high-error regions require explicit treatment when curiosity signals are used for policy optimization \cite{mavor-parker-2021}.

This study addresses the gap by evaluating an intrinsic-reward family based on feature-space impact and region-local learning progress.
The base reward combines the magnitude of latent transition change with an estimate of local model improvement.
The gated variant further disables intrinsic shaping in regions that satisfy a high-error low-progress criterion, thereby treating unproductive curiosity as a local property rather than only as a global scheduling problem.
This design follows the literature's motivation for curiosity, impact, and progress, while focusing on their interaction within a PPO-based deep reinforcement-learning setting.

The research gap is therefore not the absence of intrinsic motivation methods, but the need for a selective intrinsic reward that distinguishes useful exploration from persistent surprise.
The combination of impact, local progress, and gating clarifies how intrinsic rewards can be made more reliable under sparse rewards, noisy dynamics, and continuous-control learning.